\documentclass[10pt,a4paper]{amsart}
\usepackage[margin=19mm]{geometry}
\usepackage{amsmath,amssymb,mathtools}
\usepackage[scr=rsfs,cal=euler]{mathalfa}
\usepackage{xfrac}
\usepackage[unicode,hidelinks,bookmarks=false]{hyperref}
\newcommand{\ie}{i.e.\ }
\newcommand{\cf}{cf.\ }
\newcommand{\resp}{respectively\ }
\newcommand{\Subsection}{\subsection}
\providecommand{\nobreakdash}{\nobreak\hbox{-}}
\setlength{\emergencystretch}{2em}
\multlinegap0pt
\usepackage{amssymb}
\usepackage{bbm}
\usepackage{nicefrac}
\usepackage{enumerate}
\usepackage{tikz}
\usepackage{mathtools}
\newtheorem{theorem}{Theorem}
\newcommand{\A}{\mathcal{A}}
\newcommand{\B}{\mathcal{B}}
\renewcommand{\P}{\mathcal{P}}
\renewcommand{\a}{\alpha}
\renewcommand{\b}{\beta}
\newcommand{\brackets}[1]{\left\langle #1 \right\rangle}
\newcommand{\cLP}{\dagger(\mathsf{LP})}
\newcommand{\card}[1]{\left\lvert #1 \right\rvert}
\newcommand{\dlt}{\delta}
\newcommand{\explicitSet}[1]{\left\lbrace #1 \right\rbrace}
\newcommand{\pwmf}{\nicefrac{\mathcal{P}(\w)}{\mathrm{Fin}}}
\newcommand{\rest}{\!\restriction\!}
\newcommand{\s}{\sigma}
\newcommand{\seq}[2]{\brackets{#1 \colon #2}}
\newcommand{\set}[2]{\explicitSet{#1 \colon #2}}
\newcommand{\sub}{\subseteq}
\newcommand{\w}{\omega}
\title{A parametrized $\diamondsuit$ for the Laver property \\ and nontrivial automorphisms of $\mathcal P(\omega)/\mathrm{Fin}$}
\author{Will Brian, Alan Dow}
\date{Source: arXiv:2601.19999v1 (CC BY 4.0)}
\begin{document}
\maketitle

\noindent\emph{Source and licence.} A parametrized $\diamondsuit$ for the Laver property \\ and nontrivial automorphisms of $\mathcal P(\omega)/\mathrm{Fin}$. Version arXiv:2601.19999v1 (CC BY 4.0), distributed by arXiv under the Creative Commons Attribution 4.0 International licence, http://creativecommons.org/licenses/by/4.0/, as recorded in the arXiv licence metadata for this exact version and shown on the abstract page https://arxiv.org/abs/2601.19999v1. Full licence notice: \texttt{licenses/arxiv-2601.19999-CC-BY-4.0.txt}.\par\smallskip

\noindent\textit{Editorial scope.} Counted unit: the article's theorem thm:main (source0.tex lines 575--578). The article's complete original proof of it is delivered with it (source0.tex lines 579--710, one \texttt{proof} environment of 132 lines). No mathematics is written, repaired or re-derived: the statement and the proof are the article's own lines, in the article's own notation, and no retained line is substituted or re-flowed.  No further source-local context is needed: every cross-reference of every retained line resolves inside the two delivered spans.  Nothing was omitted from any retained span, so the package carries no omission record.  No retained line contains a citation, so the wrapper carries no bibliography and no bibliography entry.  No retained line contains a figure environment, an image-inclusion command or drawing code, so no image is delivered and the package declares no asset dependency.  Editorial formatting only, in addition: an AMS \texttt{amsart} preamble that restates the article's own declarations and macros used by the retained lines, namely the environments \texttt{theorem} on separate counters, together with the standard packages those lines need.  The retained environments print in this document as Theorem 1 in order of appearance, whereas the article prints the same items with the numbering of its own full document.  The complete native source of the article as distributed in the arXiv e-print for this exact version is bundled unchanged as \texttt{sources/193498/source0.tex}.
\begin{theorem}\label{thm:main}
Suppose $\dagger(\mathsf{LP}_V)$ holds with respect to an inner model $V$ with $\w_1^V = \w_1$.  
Then there is a bowtie of functions, and consequently, there is a nontrivial automorphism of $\pwmf$. 
\end{theorem}

\begin{proof}
To begin, let $\seq{I_n}{n \in \w \setminus 3}$ denote the sequence of finite intervals in $\w$ such that 
$|I_n| = n$ and $\min(I_3) = 0$ and $\min(I_{n+1}) = 1 + \max(I_n)$ for all $n$. 
(For convenience, many countably infinite sets in this proof will be indexed with $\w \setminus 3 = \{3,4,5,\dots\}$ rather than $\w$. This simplifies some formulas later on.) 
Our goal is to use $\cLP$ to define a strictly $\sub^*$-decreasing sequence $\seq{J_\a}{\a < \w_1}$ of subsets of $\w$ such that for every $\a$, 
\begin{itemize}
\item[$\circ$] $\card{J_\a \cap I_n} \geq 2$ for all $n \geq 3$.
\item[$\circ$] Let $S_\a = \set{n \geq 3}{\card{J_\a \cap I_n} > 2}$, and let $s_\a$ denote the unique increasing bijection $\w \setminus 3 \to S_\a$. Then $\card{J_\a \cap I_{s_\a(k)}} = k$.
\end{itemize}
Our bowtie of functions will then consist of the functions 
$f_\a$ 
with domain 
$$D_\a \,=\, \textstyle \bigcup \set{J_\a \cap I_n}{\card{J_\a \cap I_n} = 2} \,=\, \bigcup_{n \in \w \setminus S_\a}J_\a \cap I_n,$$
where $f_\a$ is defined so that it swaps the two elements of $J_\a \cap I_n$ for every $n \in \w \setminus S_\a$. 
Any strictly $\sub^*$-decreasing sequence $\seq{J_\a}{\a < \w_1}$ with these properties gives rise to a sequence of permutations $f_\a: D_\a \to D_\a$ satisfying the first two items in the definition of a bowtie of functions. 
The crux of the proof is using $\cLP$ to recursively build a sequence $\seq{J_\a}{\a < \w_1}$ that also satisfies the third item in the definition. 

Before applying $\cLP$ we must choose a sequence $\seq{h_\a}{\a < \w_1}$ of functions in $\A \cap V$. 
Before doing this, however, we must first construct the sets $S_\a$ and their enumeration functions, as described above, by a transfinite recursion that takes place in $V$. 
(While $S_\a$ is defined from $J_\a$ above, it is a key feature of the proof that the $S_\a$ are defined first, before the $J_\a$. Later the $J_\a$ are chosen in such a way that the above relation between $S_\a$ and $J_\a$ holds.) 
After constructing the $S_\a$, we will define each $h_\a$ from $S_\a$. 

For the base case of the recursion, let $S_0 = \w \setminus 3$. 
For the successor step, suppose $S_\a$ is given. 
Let $s_\a$ denote the unique increasing bijection $\w \setminus 3 \to S_\a$, define $s_{\a+1}: \w \setminus 3 \to \w$ by setting $s_{\a+1}(k) = s_\a(k2^k)$ for all $k \geq 3$, and let $S_{\a+1} = \mathrm{Image}(s_{\a+1})$. 
In other words, we define $S_{\a+1}$ to be the set containing the $(k2^k - 3)^\mathrm{th}$ element of $S_\a$ for every $k \geq 3$. Clearly $S_{\a+1}$ is an infinite subset of $S_\a$, and $S_\a \setminus S_{\a+1}$ is also infinite. 
For the limit step, suppose $\seq{S_\xi}{\xi < \a}$ is given for some limit ordinal $\a$, and that it is a strictly $\sub^*$-decreasing sequence of sets. 
Let $S_\a$ be an infinite pseudo-intersection of $\set{S_\xi}{\xi < \a}$, and let $s_\a$ denote the unique increasing bijection $\w \setminus 3 \to S_\a$. Shrinking $S_\a$ further if needed, we may (and do) assume that $s_\xi \leq^* s_\a$ for all $\xi < \a$. 

As mentioned already, this entire recursion takes place in the inner model $V$, and thus the result of the recursion, the sequences $\seq{S_\a}{\a < \w_1}$ and $\seq{s_\a}{\a < \w_1}$, are members of $V$ as well. 
For each $\a < \w_1$, define $h_\a: \w \to \w$ so that $h_\a(n) = 1$ if $n < 2^{s_{\a+1}(2)}$, and otherwise  
$$h_\a(n) = j \quad \Leftrightarrow \quad n \in \big[ s_{\a+1}(j),s_{\a+1}(j+1) \big).$$
Observe that each $h_\a \in \A$, because each $s_{\a+1}$ is strictly increasing, and that $\seq{h_\a}{\a < \w_1} \in V$, because $\seq{s_\a}{\a < \w_1} \in V$. 

Fix a function $g: \w_1 \to \mathcal S \cap V$ witnessing $\cLP$ for $\seq{h_\a}{\a < \w_1}$. 
The $J_\a$ are built by transfinite recursion using $g$ as a guide. 
But before giving this construction, let us first consider one way in which subsets of $\w$ can be coded by functions $\w \to \w$. 
Define $b: \w \to \w$ by setting $b(0)=b(1)=b(2)=0$, and $b(n) = 2^{|I_n|} = 2^n$ for $n \geq 3$. 
Working in $V$, there is for each $n \geq 3$ a bijection $\phi_n$ from $b(n) = \{0,1,\dots,2^n-1\}$ to $\P(I_n)$. 
Using these bijections, there is a natural homeomorphism $\Phi$ between 
$K_b = \set{a \in \w^\w}{a \leq b}$ and $\prod_{n \in \w}\P(I_n)$. 
Thus we may (and do) think of the members of $K_b$ as ``coding'' subsets of $\w$. Specifically, $A \sub \w$ is coded by the function $n \mapsto A \cap I_n$ in $\prod_{n \in \w}\P(I_n)$, which corresponds to a member of $K_b$ via $\Phi$, and a function $a \in K_b$ codes the set $\bigcup_{n \in \w}\Phi(a)(n) = \bigcup_{n \in \w}\phi_n\big(a(n)\big)$. 

We are now ready to describe the construction of $\seq{J_\a}{\a < \w_1}$. 
For the base step, let $J_0 = \w$. Note that $|J_0 \cap I_n| = |I_n| = n$ for all $n \geq 3$, and therefore $\set{n \geq 3}{|J_0 \cap I_n| > 2} = \w \setminus 3 = S_0$. 

For the successor step, suppose $J_\a \sub \w$ and $\set{n \geq 3}{|J_\a \cap I_n| > 2} = S_\a$. Let $s_\a$ denote the unique increasing bijection $\w \setminus 3 \to S_\a$, as above, and suppose furthermore that $|J_\a \cap I_{s_\a(k)}| = k$ for all $k \geq 3$. 
We shall describe $J_{\a+1}$ by determining the value of $J_{\a+1} \cap I_n$ for every $n \in \w \setminus 3$. 
First, if $|J_\a \cap I_n| = 2$ (i.e., if $n \notin S_\a$) then set $J_{\a+1} \cap I_n = J_\a \cap I_n$. 
Second, if $n \in S_\a \setminus S_{\a+1}$, then let $J_{\a+1} \cap I_n$ be the first two elements of $J_\a \cap I_n$. 
Finally, consider $n \in S_{\a+1}$. 
Recall that $s_{\a+1}$ denotes the unique increasing bijection $\w \setminus 3 \to S_{\a+1}$, and let $k = s_{\a+1}^{-1}(n)$. 
Recall that  $s_{\a+1}(k) = s_\a(k2^k)$, and $|J_\a \cap I_n| = |J_\a \cap I_{s_{\a+1}(k)}| = |J_\a \cap I_{s_\a(k2^k)}| = k2^k$. Our goal is to choose, from among the $k2^k$ elements of $J_\a \cap I_n$, some $k$ of them to be $J_{\a+1} \cap I_n$. 

Recall that $g(\a)$ is an $h(\a)$-slalom, meaning that $g(\a)(n)$ consists of $h(\a)(n)$ finite sequences in $\w^{n+1}$, and for each $\s \in h(\a)(n)$, we have $\s(n) \in \w$. If $\s(n) < 2^n$ then $\s(n)$ codes a subset of $I_n$ via $\phi_n$. 
Let 
$$\B_k = \set{\phi_n\big(\s(n)\big) \sub I_n}{\s \in g(\a)(n) \text{ and }\s(n) < 2^n}.$$ 
In other words, $g(\a)$ is telling us, via coding, a set of subsets of $I_n$. 
Furthermore, $h_{\a}(n) = h_{\a}(s_{\a+1}(k)) = k$, so 
$\B_k$ is a set of at most $k$ subsets of $I_n$. 
(While $g(\a)(n)$ has size exactly $k$, some of the $\s \in g(\a)(n)$ may have $\s(n) \geq 2^n$, in which case they do not code subsets of $I_n$ and do not contribute to $\B_k$.) 
These $k$ sets generate a partition of $I_n$ with size at most $2^k$. 
Since $|J_\a \cap I_n| = k2^k$, there are some $k$ members of $J_\a \cap I_n$ that are all in the same piece of the partition. 
Therefore we may (and do) choose $J_{\a+1} \cap I_n$ to be a size-$k$ subset of $J_\a \cap I_n$ such that $J_{\a+1} \cap I_n$ is either contained in or disjoint from every $B \in \B_k$. 
This concludes the successor step of the construction of the $J_\a$. 
It is clear that our recursive hypotheses are preserved: we have $J_{\a+1} \sub J_\a$, and $S_{\a+1} = \set{n \geq 3}{|J_\a \cap I_n| > 2}$, and $|J_{\a+1} \cap I_{s_{\a+1}(k)}| = k$ for all $k \geq 3$. 

For the limit step, fix a limit ordinal $\a < \w_1$, and recall that $S_\a$ is a pseudo-intersection of $\set{S_\xi}{\xi < \a}$, and that $s_\xi \leq^* s_\a$ for all $\xi < \a$. 
We will define $J_\a$ in two steps: first we define a set $J^0_\a$ by describing $J^0_\a \cap I_n$ for all $n \geq 3$ (like with $J_{\a+1}$ in the successor step), and then we shrink $J^0_\a$ to the desired set $J_\a$. 
As a recursive hypothesis (one clearly preserved at successor steps, where $J_{\a+1} \sub J_\a$), assume $\seq{J_\xi}{\xi < \a}$ is $\sub^*$-decreasing. 

Fix an increasing sequence $\seq{\xi_n}{n < \w}$ of ordinals below $\a$ with $\xi_0 = 0$ and $\sup_{n < \w}\xi_n = \a$. 
Next, fix a strictly increasing sequence $\seq{k_n}{n < \w}$ of natural numbers with $k_0 > 3$ 
such that 
\begin{itemize}
\item[$\circ$] $S_\a \setminus k_n \sub S_{\xi_{n+1}}$, and 
\item[$\circ$] if $i \geq k_n$ then $s_{\xi_{n+1}}(i) \leq s_\a(i)$, and 
\item[$\circ$] $J_{\xi_{n+1}} \!\setminus s_{\xi_n}(k_n) \sub J_{\xi_n}$. 
\end{itemize}
It is possible to find such a sequence because $S_\a \sub^* S_{\xi_{n+1}}$, and $s_{\xi_{n+1}} \leq^* s_\a$, and $J_{\xi_{n+1}} \sub^* J_{\xi_n}$ for all $n$, and because each $s_{\xi_n}$ is strictly increasing. 
Define $J^0_\a \sub \w$ by taking
\begin{align*}
J^0_\a \cap I_j &= J_0 \cap I_j = I_j \qquad \text{ if } \ j \in \big[ 3,s_0(k_0) \big), \\
J^0_\a \cap I_j &= J_{\xi_{n+1}} \cap I_j \qquad \qquad\hspace{-4.2mm} \text{ if } \ j \in \big[ s_{\xi_n}(k_n),s_{\xi_{n+1}}(k_{n+1}) \big), 
\end{align*}

For each $n \in \w$, the third requirement in our choice of $k_n, k_{n-1}, \dots, k_0$ ensures that if $j \in \big[ s_{\xi_n}(k_n),s_{\xi_{n+1}}(k_{n+1}) \big)$ then 
$$J^0_\a \cap I_j = J_{\xi_{n+1}} \cap I_j \sub J_{\xi_n} \cap I_j \sub J_{\xi_{n-1}} \cap I_j \sub \dots \sub J_{\xi_0} \cap I_j.$$ 
That is, $J^0_\a \cap I_j \sub J_{\xi_\ell} \cap I_j$ for all $\ell \leq n+1$ when $j \in \big[ s_{\xi_n}(k_n),s_{\xi_{n+1}}(k_{n+1}) \big)$. 
It follows that $J_{\xi_\ell} \sub^* J_\a^0$ for all $\ell \in \w$. 
Furthermore, if $\xi < \a$ then there is some $\ell \in \w$ with $\xi < \xi_\ell$, hence $J_\xi \supseteq^* J_{\xi_\ell} \supseteq^* J^0_\a$. Consequently, $J^0_\a \sub^* J_\xi$ for every $\xi < \a$. 

For each $n \in \w$, the first two requirements in our choice of $k_n$ ensure that if $j \in \big[ s_{\xi_n}(k_n),s_{\xi_{n+1}}(k_{n+1}) \big)$ and $j = s_\a(i)$ for some $i$, then $s_\a(i) = s_{\xi_{n+1}}(i')$ for some $i' \geq i$. (The fact that $s_\a(i) \in S_\a$ implies $s_\a(i) \in S_{\xi_{n+1}}$ by the first requirement, and $s_{\xi_{n+1}}(i) \leq s_\a(i)$ by the second requirement. $s_\a(i) \in S_{\xi_{n+1}}$ implies $s_\a(i) = s_{\xi_{n+1}}(i')$ for some $i'$, and then $i \leq i'$ because $s_{\xi_{n+1}}(i) \leq s_\a(i)$ and $s_{\xi_{n+1}}$ is strictly increasing.) 
In this case, we have $|J^0_\a \cap I_{s_\a(i)}| = |J_{\xi_{n+1}} \cap I_{\s_{\xi_{n+1}}(i')}| = i' \geq i$. 
As this is true for every $n$, it follows that if $s_\a(i) > s_{\xi_0}(k_0)$ then $|J^0_\a \cap I_{s_\a(i)}| > i$. If $s_\a(i) < s_{\xi_0}(k_0)$ then $|J^0_\a \cap I_{s_\a(i)}| = |I_{s_\a(i)}| = s_\a(i) \geq i$, as $s_\a$ is strictly increasing. This shows that $\set{n \geq 3}{|J_\a^0 \cap I_n| > 2} \supseteq S_\a$, and that 
$|J_\a^0 \cap I_{s_\a(i)}| \geq i$ for all $i$. 

By the previous paragraph, it is possible to shrink $J_\a^0$ to a set $J_\a \sub J_\a^0$ such that 
$\set{n \geq 3}{|J_\a^0 \cap I_n| > 2} = S_\a$ and 
$|J_\a^0 \cap I_{s_\a(i)}| = i$ for all $i$. 
By the paragraph before the last one, we also have $J_\a \sub J^0_\a \sub^* J_\xi$ for all $\xi < \a$. 
This completes the limit step of the recursion. 

We can now define a bowtie of functions from the sequence $\seq{J_\a}{\a < \w_1}$, as outlined at the beginning of the proof: for each $\a < \w_1$ let 
$$D_\a \,=\, \textstyle \bigcup \set{J_\a \cap I_n}{\card{J_\a \cap I_n} = 2} \,=\, \bigcup_{n \in \w \setminus S_\a}J_\a \cap I_n,$$
and let $f_\a$ be the permutation of $D_\a$ that swaps the two elements of $J_\a \cap I_n$ for every $n \in \w \setminus S_\a$. Clearly each $f_\a$ has no fixed points. 

Recall that the $S_\a$ and the $J_\a$ are strictly $\sub^*$-decreasing. 
This implies the $D_\a$ are strictly $\sub^*$-increasing. 
And it is clear that if $\a \leq \b$ then, because $J_\b \sub^* J_\a$ and $J_\a$ and $J_\b$ both meet each $I_n$ in at least two points, it is true for all but finitely many $n$ that if $|J_\a \cap I_n| = 2$ then $J_\a \cap I_n = J_\b \cap I_n$. It follows that $f_\a =^* f_\b \rest D_\a$. 

Thus $\seq{f_\a}{\a < \w_1}$ satisfies the first two requirements in the definition of a bowtie of functions, and it remains to show that it satisfies the third. 
Fix $A \sub \w$. 
Let $a$ be the function that codes $A$ as described above: i.e., $A = \bigcup_{n \in \w}\Phi(a)(n) = \bigcup_{n \in \w}\phi_n\big(a(n)\big)$. 
Recall that $a(n) < 2^{|I_n|} = 2^n$ for all $n \geq 3$, which means that $a$ is bounded by a function in $V$. 
Because $g$ witnesses $\cLP$ for $\seq{h_\a}{\a < \w_1}$, the set 
$\set{\a < \w_1}{g(\a) \text{ captures } a}$ is nonempty. 
Fix some $\dlt < \w_1$ such that $g(\dlt)$ captures $a$. 

Let $\dlt_A = \dlt+1$. 
At stage $\dlt_A$ of our construction, when we defined $J_{\dlt_A}$, 
for every $n = s_{\dlt_A}(k) \in S_{\dlt_A}$ we chose $J_{\dlt_A} \cap I_n$ to be a size-$k$ subset of $J_{\dlt} \cap I_n$ such that $J_{\dlt_A} \cap I_n$ is either contained in or disjoint from every $B \in \B_k$, where 
$$\B_k = \set{\phi_n\big(\s(n)\big) \sub I_n}{\s \in g(\dlt)(n) \text{ and }\s(n) < 2^n}.$$ 
Because $g(\dlt)$ captures $a$, we have $a(n) \in \set{\s(n)}{\s \in g(\dlt)(n)}$ for every $n$. 
Consequently, $A \cap I_n = \phi_n\big(a(n)\big) \in \B_k$ for every $n$. 
Thus $J_{\dlt_A} \cap I_n$ is either contained in or disjoint from $A \cap I_n$ for every $n \in S_\a$. 

Now fix $\a$ with $\dlt_A \leq \a < \w_1$. 
We have $D_{\dlt_A} \sub^* D_\a$, and furthermore, $D_\a \setminus D_{\dlt_A}$ is the infinite union of sets of the form $J_\a \cap I_n$ with $|J_\a \cap I_n| = 2$, where $n \in S_{\dlt_A} \!\setminus S_\a$. 
And for all but finitely many such $n$, $J_\a \cap I_n \sub J_{\dlt_A} \cap I_n$ (because $J_\a \sub^* J_{\dlt_A}$). By the previous paragraph, the sets $J_{\dlt_A} \cap I_n$, for $n \in S_{\dlt_A}$, were chosen so that $J_{\dlt_A} \cap I_n$ is either contained in or disjoint from $A$. 
Since $f_\a$ simply swaps these pairs of points, we see that in all but finitely many cases, $f_\a \setminus f_{\dlt_A}$ swaps two points in $A$ with each other or two points not in $A$ with each other. In other words, $D_\a \setminus D_{\dlt_A}$ is almost $f_\a$-invariant.
\end{proof}

\end{document}
